\section{Weekly Summary}

\begin{tcolorbox}[colback=yellow!10,colframe=black!50,title=AI-Generated Content]
\noindent The summary below was written by Claude (Anthropic) from that week's regression outputs and series descriptions. It has not been reviewed or fact-checked by a human, and should be read as an automated, informal recap -- not analysis.
\end{tcolorbox}

Welcome back to the weekly ritual where a computer flings two unrelated economic time series at each other and I pretend to be surprised by the results. This week's grab bag included Bitcoin, a discontinued San Francisco tech index, Italian central bank foreign exchange purchases, delinquency rates on real estate loans, NAFTA recession flags, U.S. exports to "NICS," a mouthful of a consumer price index, and, for continental flavor, French dwelling construction starts. Truly the theme was "no theme," which is exactly how we like it.

Monday kicked off strong with Bitcoin versus the San Francisco Tech Pulse, and the raw regression came in swinging with an R-squared of about 0.70 and a p-value so small it needs scientific notation to feel loved. Naturally I got excited, then remembered these are two things that both went up over the years, which is the correlation equivalent of noticing that both my houseplant and my nephew got taller. But here's the mildly interesting part: after detrending, it didn't fully collapse. The relationship stayed significant, just weaker, with R-squared sliding down to a humbler 0.14. So while I am legally and morally obligated to tell you this proves nothing, it's one of the rare weeks where a link survived having its shared drift surgically removed. Put a gold star on it and move on.

The rest of the week was a parade of duds, which honestly is the more honest outcome. Tuesday's Philadelphia capital expenditures diffusion index versus real user cost of money posted a p-value of 0.89, meaning "no." Wednesday's Italian foreign exchange purchases looked mildly significant against loan delinquency rates in the raw version, then evaporated the second we detrended it, p-value strolling up to 0.35 without a care in the world. That, friends, is the textbook trend-mirage this project exists to point and laugh at. Wednesday's NAFTA recession indicator versus current account balance never even bothered to look interesting, and Thursday's U.S. exports versus state-and-local government spending share flipped its coefficient sign between raw and detrended, which is the statistical version of a witness changing their story on the stand.

The week's closing weirdness came Friday, when a consumer price index for nondurables cozied up to French residential construction starts. The raw regression shrugged (p-value 0.12), but the detrended version squeaked under the significance line at 0.04, insisting on a positive relationship. Do American nondurable prices secretly whisper sweet nothings to French homebuilders? Almost certainly not, and with an R-squared of 0.08 it's a whisper at best. As always, the standing disclaimer applies with full force: these are randomly paired series, no causation is claimed, no theory is endorsed, and spurious correlation is the house specialty rather than a bug. See you next week for more numbers that mean absolutely nothing together.
